% !TeX root = ../本科微分几何笔记.tex
% 由主文档通过 \input 载入；请编译主文档。
\setcounter{chapter}{6} % 保留原书章号。
\chapter{旋转曲面、直纹面与极小曲面}
\label{chap:dg-examples}

\section{旋转曲面的统一公式}
设平面母线 $(r(s),z(s))$ 按弧长参数化，即 $r'^2+z'^2=1$，且 $r(s)>0$。绕 $z$ 轴旋转，得到
\[
 \mathbf x(s,\theta)=(r(s)\cos\theta,r(s)\sin\theta,z(s)).
\]
选择 $\nu=(-z'\cos\theta,-z'\sin\theta,r')$，直接计算得
\begin{equation}
 \mathrm I=\dd s^2+r^2\dd\theta^2,\qquad
 \mathrm{II}=(r'z''-z'r'')\dd s^2+rz'\dd\theta^2.
 \label{eq:dg-revolution-forms}
\end{equation}
坐标方向互相正交，且第二基本形式也对角化，因此母线方向和纬线方向是主方向，其主曲率为
\[
 k_s=r'z''-z'r'',\qquad k_\theta=z'/r.
\]
由单位速关系求导得 $r'r''+z'z''=0$，故
\begin{equation}
 K=-\frac{r''}{r}.
 \label{eq:dg-revolution-K}
\end{equation}
推导中不必除以 $z'$：直接使用 $(r'z''-z'r'')z'=-r''$，所以即使 $z'=0$ 也成立。条件 $r>0$ 排除了旋转轴上的坐标退化点，若要加入极点必须另查曲面的光滑性。

\begin{example}[环面的正负曲率]\label{exam:dg-torus}
对 $R>a>0$，环面参数为
\[
 \mathbf x(u,v)=((R+a\cos u)\cos v,(R+a\cos u)\sin v,a\sin u).
\]
取外法向 $\nu=(\cos u\cos v,\cos u\sin v,\sin u)$。基本形式系数为
\[
 E=a^2,\ F=0,\ G=(R+a\cos u)^2,\quad
 e=-a,\ f=0,\ g=-(R+a\cos u)\cos u.
\]
所以 $k_1=-1/a$、$k_2=-\cos u/(R+a\cos u)$，并有
\[
 K=\frac{\cos u}{a(R+a\cos u)}.
\]
外侧为正曲率，内侧为负曲率，顶部与底部的两条圆为抛物点集合。闭曲面的曲率不必处处同号。
\end{example}

\section{直纹面与可展性}
\begin{definition}[直纹面]
形如 $\mathbf x(u,v)=c(u)+v\,d(u)$、$\norm{d(u)}=1$ 的参数曲面称为\emph{直纹面（ruled surface）}，其中固定 $u$、改变 $v$ 得到的直线为母线。只在 $\mathbf x_u\times\mathbf x_v\ne0$ 的区域讨论其正则部分。
\end{definition}
由于 $\mathbf x_{vv}=0$，第二基本形式的 $g=0$，故 $K=-f^2/(EG-F^2)\le0$。计算
\[
 f=\ip{d'}\nu
 =\frac{\det(d',c'+vd',d)}{\norm{(c'+vd')\times d}}
 =\frac{\det(d',c',d)}{\norm{(c'+vd')\times d}},
\]
并利用 $EG-F^2=\norm{\mathbf x_u\times\mathbf x_v}^2$，得
\begin{equation}
 K=-\frac{\det(c',d,d')^2}{\norm{(c'+vd')\times d}^4}.
 \label{eq:dg-ruled-K}
\end{equation}
因此，直纹面正则部分满足 $K=0$ 当且仅当 $c',d,d'$ 线性相关。这也等价于法向沿母线不变：$g=f=0$ 意味着 $S\mathbf x_v=0$，即 $\nu_v=0$。满足此条件的直纹面称为\emph{可展直纹面（developable ruled surface）}。一般曲面“局部可等距展开到平面”的内蕴意义在第~\ref{chap:dg-connection} 章说明。

圆柱对应 $d$ 为常向量；圆锥对应所有母线经过一点，应去掉顶点。空间曲线 $\gamma$ 的\emph{切线曲面（tangent developable）}为 $\mathbf x(s,v)=\gamma(s)+v\mathbf t(s)$，若曲线单位速且 $\kappa>0$，则叉积大小为 $|v|\kappa$，故 $v=0$ 处不正则，而 $v\ne0$ 的局部正则部分满足 $K=0$。不能把产生奇点的参数线直接纳入正则曲面理论。

\section{极小曲面与面积的一阶变化}
\begin{definition}[极小曲面]
平均曲率处处为零的曲面称为\emph{极小曲面（minimal surface）}。条件 $H=0$ 与法向的正负选择无关。“极小”这一术语指面积的驻定条件，不保证任意大小的区域都是面积最小者。
\end{definition}
对例~\ref{exam:dg-graph-curvature} 中的图像，$H=0$ 等价于
\begin{equation}
 (1+h_v^2)h_{uu}-2h_uh_vh_{uv}+(1+h_u^2)h_{vv}=0.
 \label{eq:dg-minimal-equation}
\end{equation}
该方程还可写为 $\operatorname{div}(\grad h/\sqrt{1+\norm{\grad h}^2})=0$。当一阶导数很小时，忽略二次梯度项可得到线性近似 $h_{uu}+h_{vv}=0$；这只是小斜率近似，不是一般极小曲面方程。

\begin{proposition}[法向变分的面积第一变分]
设紧致曲面区域 $D$ 的法向为 $\nu$，$f$ 为在边界附近为零的光滑函数。对 $\mathbf x_\varepsilon=\mathbf x+\varepsilon f\nu$，在充分小 $\varepsilon$ 下保持浸入，有
\begin{equation}
 \left.\frac{\dd}{\dd\varepsilon}\right|_{0}\operatorname{Area}(D_\varepsilon)
 =-2\int_D fH\dd A.
 \label{eq:dg-area-variation}
\end{equation}
因此，所有这类法向变分下的面积一阶变化为零，当且仅当 $H=0$。
\end{proposition}
\begin{proof}
记变分度量为 $g_{ij}(\varepsilon)$。由于 $\ip\nu{\mathbf x_i}=0$，有 $\dot g_{ij}(0)=f\ip{\nu_i}{\mathbf x_j}+f\ip{\mathbf x_i}{\nu_j}=-2fb_{ij}$。行列式微分公式给出
\[
 \left.\frac{\dd}{\dd\varepsilon}\right|_0\sqrt{\det\mathbf g(\varepsilon)}
 =\tfrac12\sqrt{\det\mathbf g}\,\tr(\mathbf g^{-1}\dot{\mathbf g})
 =-2fH\sqrt{\det\mathbf g}.
\]
紧致性支持积分号下求导，得到公式。若对所有紧支撑 $f$ 积分都为零而某点 $H\ne0$，可在 $H$ 保持同号的小邻域选同号非负测试函数，使积分不为零，矛盾。
\end{proof}

\begin{example}[悬链面与正螺旋面]\label{exam:dg-minimal-models}
悬链面（catenoid）的参数为 $\mathbf x(u,v)=(a\cosh v\cos u,a\cosh v\sin u,av)$，$a>0$。取 $\nu=(\cos u/\cosh v,\sin u/\cosh v,-\tanh v)$，则
\[
 E=G=a^2\cosh^2v,\quad F=0,\qquad e=-a,\ f=0,\ g=a.
\]
故 $H=0$、$K=-1/(a^2\cosh^4v)$。正螺旋面（helicoid）取 $\mathbf y(u,v)=(u\cos v,u\sin v,av)$，其度量为 $\dd u^2+(u^2+a^2)\dd v^2$；以叉积方向作法向，第二基本形式系数为 $e=g=0$、$f=-a/\sqrt{u^2+a^2}$，同样有 $H=0$，且 $K=-a^2/(u^2+a^2)^2$。它还是直纹面，但不是可展面，说明“由直线生成”不等于“可以等距展开”。
\end{example}

\section{习题}
\begin{exercise}\label{ex:dg-examples-cone}
圆锥去掉顶点后取 $\mathbf x(r,\theta)=(r\cos\theta,r\sin\theta,cr)$，$r>0$。求度量，并构造到平面扇形的局部等距坐标。
\end{exercise}
\begin{exercise}\label{ex:dg-examples-torus}
用例~\ref{exam:dg-torus} 的公式计算整个环面的高斯曲率积分，解释正负曲率如何抵消。
\end{exercise}
\begin{exercise}\label{ex:dg-examples-helicoid}
验证例~\ref{exam:dg-minimal-models} 中正螺旋面的基本形式与 $K,H$。
\end{exercise}
\begin{exercise}\label{ex:dg-examples-sphere-variation}
用球面积 $4\pi R^2$ 直接计算外法向常速变分 $R\mapsto R+\varepsilon$ 的面积导数，并与式~\eqref{eq:dg-area-variation} 核对符号。
\end{exercise}

